Computer-use agents often fail by announcing that a task is finished when it is not. We diagnose it in Holo-3.1-4B on real OSWorld desktops and test whether editing activations or weights, without training, removes it. The agent scores \oswMeanScorePct\% on \oswNScored{} diagnostic tasks; a false ``done'' is \oswFsPct\% of its failures, and re-sampled from the same contexts it answers ``done'' on false and true completions alike. With an exact replay pipeline and a pre-registered gate, we find three traps that would each have produced a positive result: a false-success direction that adds at most \ofPosGain{} AUROC to position, run-to-run noise under which only \rtmFsAgainPct\% of false successes recur, and a loop direction in a smaller model that beats position yet fails every causal test. No steering arm or weight edit shows a large effect in pilot screens, which flag only a fall from 8/8 to \gateMaxAnswers/8 false answers and hold a single feasible false-success context, and no arm run online beats the run-to-run noise of the original. The tool call leans on a thought that already declares success: on \paN{} feasible contexts, neutral text in its place removes about 80\% of the \texttt{answer} margin, yet \texttt{answer} stays preferred in \paNeutralPosAll; the note reinforces the conclusion rather than creating it. Of the edits that move the call, on \pwOrigTasks{} feasible false successes, a model merge lowers true completions as much as false ones, a lever without a discriminator, and a closed-form edit lowers false ones more, with a selectivity interval that includes zero. An external verifier that checks each requirement on the final screen does separate them offline: it detects \gbClaudeDet\% of false successes and rejects \gbClaudeRej\% of true ones, against \gbHoloDet\% and \gbHoloRej\% for the agent judging itself; its effect on the OSWorld score is not yet measured.
